\section{Introduction}

Research code for driving agents often falls into one of two extremes: visually compelling demos with weak experimental structure, or benchmark-oriented systems that are difficult to run, inspect, and extend. Large platforms such as CARLA and MetaDrive have been influential for realism and scale, while procedural reinforcement-learning benchmarks such as Procgen and CoinRun have helped foreground generalization as a first-class evaluation target \cite{carla2017, metadrive2021, procgen2019, coinrun2018}. EvoDrive Lab is intended to occupy a different point in that design space: a lightweight, Docker-first benchmark where multiple agents can be trained, replayed, and compared on procedurally generated 2D tracks without large infrastructure overhead.

The central question of this paper is whether different training methods generalize similarly to unseen procedural tracks. A method can appear effective during optimization yet still degrade on held-out environments. EvoDrive Lab makes that gap visible and measurable by pairing shared train, validation, and test suites with replay-oriented qualitative analysis. The benchmark is small enough to run on commodity hardware, but structured enough to support repeatable comparisons and paper-ready reporting.

This paper does not claim state-of-the-art autonomous driving performance. Instead, it contributes a compact, reproducible benchmark and an empirical comparison of a custom genetic algorithm, NEAT \cite{stanley2002}, and a lightweight PPO-style baseline inspired by proximal policy optimization \cite{schulman2017}. The emphasis is on comparative behavior inside one controlled benchmark rather than broad conclusions about autonomous driving.

More concretely, the paper makes three contributions:
\begin{itemize}[leftmargin=1.5em]
    \item a small procedural 2D driving benchmark that links quantitative evaluation with qualitative replay artifacts;
    \item a reproducible comparison of neuroevolution and lightweight policy optimization under a shared train/validation/test protocol;
    \item two compact ablations that test how sensor richness and track length affect generalization.
\end{itemize}
